\documentclass{article}
\usepackage{graphicx} % Required for inserting images
\usepackage{amsmath}
\usepackage{amsfonts}
\usepackage{booktabs}
\usepackage{multirow}
\usepackage{float}
\usepackage{geometry}
\geometry{a4paper, margin=1in}

\title{Deep Learning Homework 1 \\ Image Classification with PyTorch on CIFAR-10}
\author{Dai Xunlian (225040015)}
\date{October 2025}

\begin{document}

\maketitle

\section{Introduction}

This report presents a comprehensive study on image classification using PyTorch on the CIFAR-10 dataset. The objective is to implement and experiment with various deep learning techniques to improve classification accuracy. We systematically evaluate different factors including network architecture modifications, optimization strategies, and regularization techniques.

The CIFAR-10 dataset consists of 60,000 32×32 color images in 10 classes, with 50,000 training images and 10,000 test images. Our experiments focus on building upon a base CNN architecture and incrementally adding improvements to achieve higher test accuracy.

\section{Experimental Setup}

\subsection{Dataset and Preprocessing}

The CIFAR-10 dataset was loaded using PyTorch's built-in datasets module. Data preprocessing involved normalizing all channels with mean and standard deviation of 0.5, converting images to PyTorch tensors, and using a batch size of 64 for both training and testing phases. This preprocessing approach ensures consistent data distribution and efficient batch processing during model training.

\subsection{Base Architecture}

Our base CNN architecture employs a carefully designed structure that balances complexity with computational efficiency. The network begins with 5 convolutional layers featuring channel dimensions of [3, 128, 256, 512, 512, 256], progressively increasing the feature map depth to capture hierarchical patterns in the images. Max pooling layers are strategically placed after layers 1 and 4 to reduce spatial dimensions while preserving important features. To prevent overfitting, dropout layers with a rate of 0.3 are applied after pooling operations. The fully connected section consists of four layers with dimensions [flatten, 512, 256, 128, 10], incorporating ReLU activation functions and dropout layers with a 0.5 dropout rate to enhance generalization capability.

\subsection{Training Configuration}

All experiments were conducted using consistent training parameters to ensure fair comparison across different architectural modifications. The learning rate was set to 0.001, providing a balanced approach between convergence speed and stability. Training was limited to a maximum of 200 epochs, with early stopping implemented using a patience value of 10 epochs to prevent overfitting and reduce computational overhead. The cross-entropy loss function was employed for all experiments, as it is well-suited for multi-class classification tasks. The primary optimizer was SGD with momentum set to 0.9, unless specifically testing alternative optimization strategies.

\section{Experimental Factors and Results}

We conducted 7 different experiments to evaluate various factors that could improve classification accuracy:

\subsection{Experiment 1: Base Model}
\textbf{Configuration:} Standard CNN architecture as described above.
\textbf{Result:} 84.10\% test accuracy
\textbf{Analysis:} This serves as our baseline for comparison with other configurations.

\subsection{Experiment 2: Increased Hidden Dimensions}
\textbf{Configuration:} Doubled the hidden dimensions in all layers (channels: [3, 256, 512, 1024, 1024, 512]).
\textbf{Result:} 84.42\% test accuracy
\textbf{Analysis:} Slight improvement (+0.32\%) suggests that increased capacity helps, but the improvement is marginal, indicating potential overfitting concerns.

\subsection{Experiment 3: Increased Layer Depth}
\textbf{Configuration:} Added more layers to create a deeper network (7 layers instead of 5).
\textbf{Result:} 82.60\% test accuracy
\textbf{Analysis:} Performance decreased (-1.50\%), likely due to vanishing gradient problems and increased training difficulty without proper residual connections.

\subsection{Experiment 4: Average Pooling}
\textbf{Configuration:} Replaced max pooling with average pooling.
\textbf{Result:} 82.27\% test accuracy
\textbf{Analysis:} Performance decreased (-1.83\%), suggesting that max pooling is more effective for preserving important features in this dataset.

\subsection{Experiment 5: Residual Connections (ResNet)}
\textbf{Configuration:} Implemented residual blocks with skip connections.
\textbf{Result:} 89.42\% test accuracy
\textbf{Analysis:} Significant improvement (+5.32\%), demonstrating the effectiveness of residual connections in solving vanishing gradient problems and enabling deeper networks.

\subsection{Experiment 6: L2 Regularization}
\textbf{Configuration:} Added L2 regularization with λ = 0.001.
\textbf{Result:} 84.13\% test accuracy
\textbf{Analysis:} Minimal improvement (+0.03\%), suggesting that dropout already provides sufficient regularization for this dataset size.

\subsection{Experiment 7: Adam Optimizer}
\textbf{Configuration:} Replaced SGD with Adam optimizer.
\textbf{Result:} 80.00\% test accuracy
\textbf{Analysis:} Performance decreased (-4.10\%), indicating that SGD with momentum works better for this specific architecture and dataset.

\section{Results Summary}

\begin{table}[H]
\centering
\begin{tabular}{@{}lcccc@{}}
\toprule
\textbf{Experiment} & \textbf{Final Accuracy} & \textbf{Best Accuracy} & \textbf{Training Time (s)} & \textbf{Improvement} \\
\midrule
Base Model & 84.10\% & 84.10\% & 787.82 & - \\
Increased Hidden Dim & 84.42\% & 84.42\% & 1357.04 & +0.32\% \\
Increased Layer Depth & 82.60\% & 82.60\% & 1055.41 & -1.50\% \\
Average Pooling & 82.27\% & 82.27\% & 931.65 & -1.83\% \\
ResNet & \textbf{89.42\%} & \textbf{89.42\%} & 851.07 & \textbf{+5.32\%} \\
L2 Regularization & 84.13\% & 84.13\% & 895.12 & +0.03\% \\
Adam Optimizer & 80.00\% & 80.00\% & 562.92 & -4.10\% \\
\bottomrule
\end{tabular}
\caption{Experimental Results Summary}
\label{tab:results}
\end{table}

\section{Analysis and Discussion}

\subsection{Key Findings}

1. \textbf{Residual Connections are Most Effective:} The ResNet configuration achieved the highest accuracy of 89.42\%, demonstrating the importance of skip connections in deep networks.

2. \textbf{Optimizer Choice Matters:} Adam optimizer performed worse than SGD with momentum, suggesting that the adaptive learning rate might be too aggressive for this dataset.

3. \textbf{Pooling Strategy Impact:} Max pooling outperformed average pooling, indicating that preserving maximum activation values is more important for feature extraction.

4. \textbf{Regularization Effects:} L2 regularization provided minimal benefit, likely because dropout already provides sufficient regularization.

5. \textbf{Network Depth vs. Width:} Increasing width slightly improved performance, while increasing depth without residual connections decreased performance.

\subsection{Architecture Insights}

The success of the ResNet configuration can be attributed to several key architectural advantages. Skip connections enable better gradient propagation during backpropagation, addressing the vanishing gradient problem that often plagues deeper networks. Residual blocks allow the network to learn identity mappings when needed, providing flexibility in feature representation. Additionally, residual connections provide alternative paths for information flow, enhancing training stability and enabling the network to effectively utilize its increased depth.

\subsection{Training Efficiency}

The ResNet configuration not only achieved the highest accuracy but also had reasonable training time (851.07 seconds), making it the most efficient solution among the successful configurations.

\section{Conclusion}

This study demonstrates the importance of architectural choices in deep learning for image classification. The key findings reveal that residual connections (ResNet) provide the most significant improvement in accuracy, achieving 89.42\% compared to the baseline 84.10\%. SGD with momentum outperforms Adam optimizer for this specific task, suggesting that adaptive learning rates may be too aggressive for this dataset. Max pooling proves more effective than average pooling for feature extraction, while additional regularization beyond dropout provides minimal benefits. Network width improvements are more effective than depth increases without residual connections, highlighting the importance of proper architectural design for deeper networks.

The best configuration achieved 89.42\% test accuracy on CIFAR-10, representing a substantial improvement over the baseline model. This demonstrates the effectiveness of systematic experimentation and incremental improvements in deep learning model development.

\section{Future Work}

Potential areas for further improvement include data augmentation techniques such as rotation, flipping, and color jittering to increase dataset diversity and improve generalization. Attention mechanisms could be incorporated to focus on relevant image regions, potentially enhancing feature extraction capabilities. Learning rate scheduling strategies might optimize training dynamics and convergence behavior. Ensemble methods could combine multiple models to achieve even higher accuracy, while transfer learning from pre-trained models on larger datasets might provide additional performance gains.

\end{document}
